% 附录：代码清单
\section{算法与实现说明}
\noindent 本附录聚焦“问题分析所用的核心算法与理论”，不再给出出图脚本流程。每章最多选取两条关键算法，并附“Python实现阐释”（文字化说明数据结构、步骤与接口），以保证专业性与可复现性。

% ==================== 问题一：确定性最优化 ====================
\subsection{切换成本的邻接线性化与增量割生成}
为避免使用不稳定的大 $M$ 常数，对相邻季作物切换采用邻接线性化，并在“违反检测—补充割”的循环中增量收紧。
\begin{algorithm}[H]
  \caption{切换成本的邻接线性化与增量割生成}
  \label{lst:switch-linearization}
  % 兼容旧引用：将“收益曲线脚本”标签指向本算法
  \label{lst:revenue-curve}
  \begin{algorithmic}[1]
    \Require 基础 MIP 模型（含 $x,z$ 与唯一选择约束），潜在切换对集合 $\mathcal{E} = \{(k\to k')\mid k\ne k'\}$
    \Ensure 收敛到无违反的可行解与切换计提变量 $w_{t,s,i,k\to k'}$
    \State 初始化：不显式生成全部 $w$ 与割，仅保留基础模型；
    \Repeat
      \State 求解当前松弛模型，得解 $(x,z)$；
      \State 按地块—季扫描相邻季 $(s,s^+)$，对每个 $i$ 与作物对 $(k\to k')$ 检测是否存在 $z_{t,s,i,k}=z_{t,s^+,i,k'}=1$；
      \State 若发现违反（未计提切换），则为对应索引生成 $w_{t,s,i,k\to k'}$ 并加入邻接不等式：\\
            $w_{t,s,i,k\to k'} \ge z_{t,s,i,k} + z_{t,s^+,i,k'} - 1$，且 $w\le z$ 的两个上界；
      \State 将 $\sum \kappa_{k\to k'}\,w_{t,s,i,k\to k'}$ 计入年度净收益；
    \Until 无新的违反被发现
    \State 返回收敛解与完整切换成本项。
  \end{algorithmic}
\end{algorithm}
\paragraph{Python实现阐释}
数据以稀疏索引存储，切换计提变量按需生成。在“求解—扫描—加约束”的循环中，先识别相邻季的潜在切换，再为对应索引追加邻接不等式并计入成本，直至不再出现新的违反。为保证稳定性可配合热启动，并记录迭代收敛过程以便复核。

% 代码清单：切换成本邻接线性化与增量割生成（含函数签名的核心片段）
\newpage
\begin{center}
\begin{minipage}[t][\dimexpr\textheight-1em\relax]{0.98\linewidth}
\lstinputlisting[
  style=appendixpy,
  firstline=11,
  lastline=63,
  caption={Python 实现：切换成本邻接线性化与增量割生成（核心）},
  label={lst:py-switch-linearization}
]{python/model1/switch_linearization.py}
\end{minipage}
\end{center}

\subsection{轮作约束的窗口化线性近似与冲突窗口检测}
对最小轮作间隔 $L_k$ 进行滑动窗口近似，结合违反检测的增量割生成，避免一次性生成指数级不等式。
\begin{algorithm}[H]
  \caption{轮作约束的窗口化线性近似与冲突窗口检测}
  \label{lst:rotation-cuts}
  % 兼容旧引用：将“堆叠面积脚本”标签指向本算法
  \label{lst:stack-area}
  \begin{algorithmic}[1]
    \Require 作物集合 $\mathcal{K}$、轮作间隔 $L_k$、季序运算 $s\oplus\Delta$
    \Ensure 满足“任意长度 $L_k$ 窗口内至多一次种植”的解
    \State 预生成窗口模板：$W_{t,s,i,k}=\{(t, s\oplus\Delta, i, k)\}_{\Delta=0}^{L_k-1}$；
    \Repeat
      \State 求解当前模型得 $z$；
      \State 对每个 $(t,s,i,k)$ 检查窗口和 $\sum\limits_{\Delta=0}^{L_k-1} z_{t, s\oplus\Delta, i, k}$；
      \State 若大于 1，向模型加入对应的窗口不等式实例；
      \State // 实现建议：预先缓存 $W$ 并用向量化索引聚合窗口和；
      \State // 将同一 $(t,i,k)$ 下的所有违反窗口批量打包写入约束，减少回调；
      \State // 对存在周期性的季序 $s$，以 $s\oplus\Delta$ 的环状移位生成窗口；
      \State // 记录“新增窗口/约束”计数，用于收敛与复杂度审计；
      \State // 可选：自适应放宽/收紧 $L_k$ 做热身，或按作物分组优先修复高违反率类别；
    \Until 所有窗口均满足上界 1
    \State 返回满足轮作间隔的解。
  \end{algorithmic}
\end{algorithm}
\paragraph{Python实现阐释}
实现包含“窗口生成器”“违反扫描器”“批量约束注入器”三组件。窗口以向量化索引快速枚举；违反窗口打包后一次性追加，有效减少求解器回调次数。

具体实现要点如下：
- 数据结构：索引采用压缩整型并以结构化数组/字典映射到变量位置，窗口模板可预先缓存以复用；
- 违反检测：通过卷积式滑窗或向量化求和高效计算窗口内计数；
- 约束注入：把同类窗口打包为稀疏矩阵块，一次性追加到求解器，显著降低回调次数；
- 可复现性：固定随机种子与索引顺序，记录每轮迭代新增的窗口与约束数量。

\newpage

% 代码清单：轮作窗口化约束与冲突检测（含函数签名的核心片段）
\begin{center}
\begin{minipage}[t][\dimexpr\textheight-1em\relax]{0.98\linewidth}
\lstinputlisting[
  style=appendixpy,
  firstline=24,
  lastline=72,
  caption={Python 实现：轮作窗口化约束与冲突检测（核心）},
  label={lst:py-rotation-cuts},
  stringstyle=\color{pyComment}\itshape,
  literate=*{∈}{{$\in$}}1 {⊕}{{$\oplus$}}1 {Δ}{{$\Delta$}}1 {Σ}{{$\Sigma$}}1 {≤}{{$\le$}}1 {≥}{{$\ge$}}1
]{python/model1/rotation_cuts.py}
\end{minipage}
\end{center}

\clearpage

% ==================== 问题二：鲁棒优化 ====================
\subsection{Bertsimas--Sim 预算不确定集的鲁棒等价}
对含不确定系数的线性约束采用 Bertsimas--Sim 预算集，将最坏情形内化为线性对偶约束，控制强度由预算参数 $\Gamma$ 调节。
\begin{algorithm}[H]
  \caption{Bertsimas--Sim 预算不确定集的鲁棒等价}
  \label{lst:bs-robust-counterpart}
  \begin{algorithmic}[1]
    \Require 名义系数 $a^0$, 扰动幅度 $d$, 右端 $b$, 预算 $\Gamma$
    \Ensure 对每条含不确定项的约束生成鲁棒等价的线性约束集
    \State 规范化元数据：校验索引对齐，定义 $\mathcal{J}=\{j\mid d_j>0\}$；
    \State 引入扩展变量：$\theta\ge0$ 与 $\{s_j\}_{j\in\mathcal{J}}\ge0$；
    \State 生成主不等式：$a^0\!\cdot x + \sum_{j\in\mathcal{J}} d_j s_j + \Gamma\,\theta \le b$；
    \State 生成上界不等式：$0\le s_j\le\theta\ \ (\forall j\in\mathcal{J})$；
    \State 边界核验：核对 $\Gamma\in\{0,|\mathcal{J}|\}$ 时与名义/全偏差的一致性；
    \State 命名与日志：统一前缀命名新增变量/约束，并记录 $\Gamma$ 与增量计数以便审计；
    \State 数值尺度：对 $a^0,d$ 做量纲尺度化以提升数值稳定性（可选）；
    \State 批处理注入并记录增量统计（变量与不等式各为 $1+|\mathcal{J}|$）；
    \State 返回扩展后的线性模型与统计。
  \end{algorithmic}
\end{algorithm}
\paragraph{Python实现阐释}
提供“鲁棒化器”函数：输入名义约束与不确定元数据，输出扩展变量与约束集合。接口支持批量处理，并记录 $(\Gamma,\text{变量/约束增量})$ 以便规模评估与复现。

实现细节与审计要点：
- 元数据规范：给出 $a^0,d,\Gamma$ 并确保索引一致；
- 模板化展开：按模板引入 $\theta, s$ 与两类线性不等式；
- 规模与审计：按约束族批处理并记录增量统计，支持复现与回放。
 - 边界核验：在小例上验证 $\Gamma=0,|\mathcal{J}|$ 的极端情形。

\newpage

% 代码清单：BS 预算集鲁棒等价（含函数签名的核心片段）
\begin{center}
\begin{minipage}[t][\dimexpr\textheight-1em\relax]{0.98\linewidth}
\lstinputlisting[
  style=appendixpy,
  firstline=29,
  lastline=79,
  caption={Python 实现：BS 预算集鲁棒等价生成器（核心）},
  label={lst:py-bs-robust}
]{python/model2/robust_bs.py}
\end{minipage}
\end{center}

\clearpage

\subsection{\texorpdfstring{$\Gamma$}{Gamma} 网格与帕累托过滤}
在 $\Gamma$ 网格上评估鲁棒解的收益与风险指标，阈值初筛后进行帕累托过滤，输出推荐的 $\Gamma^*$。
\begin{algorithm}[H]
  \caption{\texorpdfstring{$\Gamma$}{Gamma} 网格与帕累托过滤}
  \label{lst:gamma-grid}
  \begin{algorithmic}[1]
    \Require 网格 $\mathcal{G}$，阈值 $(X, Y\%)$；
    \Ensure 推荐 $\Gamma^*$ 及其对应解
    \For{$\Gamma\in\mathcal{G}$}
      \State 求解鲁棒模型得 $x^\Gamma$，计算 $r^{\min}, r^{\mathrm{nom}}, \widehat r^{\,\mathrm{avg}}$, ViolRate, AvgGap, Switches, PeakResource；
      \State // 记录求解时间与收敛状态，过滤异常点；
    \EndFor
    \State 候选集 $\mathcal{S} \gets \{\Gamma:\; r^{\min}\ge X,\; \text{ViolRate}\le Y\%\}$；
    \State 在 $\mathcal{S}$ 内做帕累托过滤：支配准则为“$r^{\min}$ 高、切换少、峰值资源低、且 $r^{\mathrm{nom}}$ 不低于阈值”；
    \State // 次序规则：先按 $r^{\min}$ 降序，其次切换次数、再峰值资源；
    \State 返回 $\Gamma^*$：若多解并列，选切换更少者。
  \end{algorithmic}
\end{algorithm}
\paragraph{Python实现阐释}
实现包含“评估器（解与指标计算）”与“过滤器（支配判定）”。支持并行评估与固定随机种子，输出含参数—指标—推荐的实验记录，便于复现与审计。

实现补充说明：
- 指标定义：$r^{\min}$（最小收益）、$r^{\mathrm{nom}}$（名义收益）、ViolRate（违反率）、AvgGap（解平均相对差距）、Switches（切换次数）、PeakResource（峰值资源占用）；
- 评估流程：在网格 $\mathcal G$ 上并行求解，归档各项指标；
- 帕累托过滤：先阈值初筛，再按“高 $r^{\min}$/低 ViolRate/低 PeakResource/不低于阈值的 $r^{\mathrm{nom}}$”进行支配判定；
- 复现性：固定随机种子与评估顺序，保存完整实验表格与推荐结果。

\newpage

% 代码清单：Gamma 网格评估与帕累托过滤（含函数签名的核心片段）
\begin{center}
\begin{minipage}[t][\dimexpr\textheight-1em\relax]{0.98\linewidth}
\lstinputlisting[
  style=appendixpy,
  firstline=5,
  lastline=60,
  caption={Python 实现：Gamma 网格评估与帕累托过滤（核心）},
  label={lst:py-gamma-grid}
]{python/model2/gamma_pareto.py}
\end{minipage}
\end{center}

\clearpage

% ==================== 问题三：价格反馈与滚动优化 ====================
\subsection{相关一致化的情景生成与尾部加密}
滚动优化依赖高质量的情景集。为在有限样本下兼顾相关结构与尾部覆盖，我们采用“PCA 对齐 + 拉丁超立方（LHS）”生成基础样本，并沿指定的尾部方向（如最小利润方向）进行尾部加密。该方法与模型三第~\ref{sec:model3} 节“相关结构与采样”一致，有助于稳定估计 CVaR。
\begin{algorithm}[H]
  \caption{相关一致化的情景生成（PCA + LHS）与尾部加密}
  \label{lst:m3-scenario-pca}
  \begin{algorithmic}[1]
    \Require 均值向量 $\mu\in\mathbb{R}^d$，协方差 $\Sigma\in\mathbb{R}^{d\times d}$，主成分数 $m$，基础样本数 $N$；可选：尾部方向 $v\in\mathbb{R}^d$ 与额外尾部样本数 $E$
    \Ensure 融合相关结构且覆盖尾部的情景样本 $X$ 与其权重 $w$
    \State 对 $\Sigma$ 做特征分解：$\Sigma=V\Lambda V^\top$，取前 $m$ 个主成分 $V_m$ 与特征值对角 $\Lambda_m$
    \State 生成 LHS 样本 $U\in(0,1)^{N\times m}$，做逐维随机置换；经标准正态逆 CDF 得 $Z=\Phi^{-1}(U)$
    \State 基础样本：$X_{\text{base}}=\mu + Z\,(V_m\Lambda_m^{1/2})^\top$
    \If{$E>0$ 且 $\|v\|>0$}
      \State 投影并单位化 $v_m= V_m^\top v/\|V_m^\top v\|$；再生成 $E$ 个 LHS 样本并取 $Z_{\text{tail}}=\Phi^{-1}(U_{\text{tail}})$
      \State 设偏置系数 $\text{scale}>0$，令 $\tilde Z=Z_{\text{tail}}-\text{scale}\cdot v_m$，得到尾部样本 $X_{\text{tail}}=\mu + \tilde Z\,(V_m\Lambda_m^{1/2})^\top$
      \State 合并 $X=[X_{\text{base}}; X_{\text{tail}}]$
    \Else
      \State $X= X_{\text{base}}$
    \EndIf
    \State 设权重 $w$ 为等权并归一化；返回 $(X,w)$
  \end{algorithmic}
\end{algorithm}
\paragraph{Python实现阐释}
本实现采用 PCA 对齐相关结构，并结合 LHS 与正态逆变换生成样本；可沿给定方向实施左尾加密以增强尾部覆盖。实现仅依赖 \texttt{numpy}，正态逆变换采用误差函数反变换思路，便于复现与审计。

% 代码清单：模型三—相关一致化的情景生成与尾部加密（核心片段）
\newpage
\begin{center}
\begin{minipage}[t][\dimexpr\textheight-1em\relax]{0.98\linewidth}
\lstinputlisting[
  style=appendixpy,
  firstline=62,
  lastline=113,
  caption={Python 实现：相关一致化的情景生成与尾部加密（核心）},
  label={lst:py-m3-scenario}
]{python/model3/scenario_generation.py}
\end{minipage}
\end{center}

\subsection{含 CVaR 的年度滚动优化（简化演示）}
为贴合模型三的“更新—生成—优化—执行—评估”闭环，这里给出一维供给决策 $q$ 的教学版实现：在年度网格上，以“期望利润$-\,\lambda\cdot\text{CVaR}@\alpha$”为目标在 $q$ 栅格中择优；随后抽取一次真实结果作为反馈，轻微更新名义参数并进入下一年。该实现对应正文中的算法~\ref{alg:rolling-cvar} 思路。
\begin{algorithm}[H]
  \caption{滚动优化 + CVaR 的年度决策流程（简化演示）}
  \label{lst:m3-rolling-cvar}
  \begin{algorithmic}[1]
    \Require 初始名义参数 $(a_0,b_0,C_0,c_0,F_0)$，年数 $T$，供给栅格 $\mathcal Q$，风险参数 $(\alpha,\lambda)$
    \For{$t=1,\dots,T$}
      \State 围绕名义参数抽取 $K$ 条情景 $\{(a,b,C,c,F)\}$（高斯扰动或由上节生成）
      \For{$q\in\mathcal Q$}
        \State 计算各情景利润 $\Pi_s(q)$ 与损失 $\ell_s(q)=-\Pi_s(q)$
        \State 估计 $\mathbb{E}[\Pi(q)]$、$\text{CVaR}_\alpha(\ell(q))$ 与 $5\%$ 分位；目标值 $J(q)=\mathbb{E}[\Pi(q)]-\lambda\,\text{CVaR}_\alpha(\ell(q))$
      \EndFor
      \State 取 $q^{(t)}=\arg\max\_{q\in\mathcal Q} J(q)$；在一条保留情景上“实现”并记录利润
      \State 依据反馈对 $(a_0,C_0)$ 做小幅漂移更新，进入下一年
    \EndFor
    \State 输出序列 $\{q^{(t)}\}$ 与对应的实现利润序列
  \end{algorithmic}
\end{algorithm}
\paragraph{Python实现阐释}
本实现包含样本 CVaR 的经验估计、批量评估过程与教学版的年度滚动主循环。价格—供给采用线性反馈并含容量上限，计算过程向量化，便于扩展到 MIP 框架；主要依赖 \texttt{numpy}。

% 代码清单：模型三—滚动优化 + CVaR（简化演示，核心片段）
\newpage
\begin{center}
\begin{minipage}[t][\dimexpr\textheight-1em\relax]{0.98\linewidth}
\lstinputlisting[
  style=appendixpy,
  firstline=91,
  lastline=139,
  caption={Python 实现：滚动优化 + CVaR（核心片段）},
  label={lst:py-m3-rolling-cvar}
]{python/model3/rolling_cvar.py}
\end{minipage}
\end{center}
